\documentclass[UTF8, 12pt, fontset=fandol, oneside]{ctexart}
\usepackage{researchnotes}

\graphicspath{{math/机器学习/}}

\title{第 10 章\quad 隐马尔可夫模型}
\author{}
\date{}

\begin{document}

\maketitle

隐马尔可夫模型（hidden Markov model, HMM）是可用于标注问题的统计学习模型，描述由隐藏的马尔可夫链随机生成观测序列的过程，属于生成模型。本章首先介绍隐马尔可夫模型的基本概念，然后分别叙述隐马尔可夫模型的概率计算算法、学习算法以及预测算法。隐马尔可夫模型在语音识别、自然语言处理、生物信息、模式识别等领域有着广泛的应用。

\section*{10.1\quad 隐马尔可夫模型的基本概念}

\subsection*{10.1.1\quad 隐马尔可夫模型的定义}

\noindent\textbf{定义 10.1（隐马尔可夫模型）}\quad 隐马尔可夫模型是关于时序的概率模型，描述由一个隐藏的马尔可夫链随机生成不可观测的状态随机序列，再由各个状态生成一个观测从而产生观测随机序列的过程。隐藏的马尔可夫链随机生成的状态的序列，称为状态序列（state sequence）；每个状态生成一个观测，而由此产生的观测的随机序列，称为观测序列（observation sequence）。序列的每一个位置又可以看作是一个时刻。

隐马尔可夫模型由初始概率分布、状态转移概率分布以及观测概率分布确定。隐马尔可夫模型的形式定义如下：

设 $Q$ 是所有可能的状态的集合，$V$ 是所有可能的观测的集合：
\[
  Q=\{q_1,q_2,\cdots,q_N\},\qquad
  V=\{v_1,v_2,\cdots,v_M\}.
\]
其中，$N$ 是可能的状态数，$M$ 是可能的观测数。

$I$ 是长度为 $T$ 的状态序列，$O$ 是对应的观测序列：
\[
  I=(i_1,i_2,\cdots,i_T),\qquad
  O=(o_1,o_2,\cdots,o_T).
\]

$A$ 是状态转移概率矩阵：
\begin{equation}
  A=[a_{ij}]_{N\times N}
  \tag{10.1}
\end{equation}
其中，
\begin{equation}
  a_{ij}=P(i_{t+1}=q_j\mid i_t=q_i),\qquad
  i=1,2,\cdots,N;\quad j=1,2,\cdots,N
  \tag{10.2}
\end{equation}
是在时刻 $t$ 处于状态 $q_i$ 的条件下在时刻 $t+1$ 转移到状态 $q_j$ 的概率。

$B$ 是观测概率矩阵：
\begin{equation}
  B=[b_j(k)]_{N\times M}
  \tag{10.3}
\end{equation}
其中，
\begin{equation}
  b_j(k)=P(o_t=v_k\mid i_t=q_j),\qquad
  k=1,2,\cdots,M;\quad j=1,2,\cdots,N
  \tag{10.4}
\end{equation}
是在时刻 $t$ 处于状态 $q_j$ 的条件下生成观测 $v_k$ 的概率。

$\pi$ 是初始状态概率向量：
\begin{equation}
  \pi=(\pi_i)
  \tag{10.5}
\end{equation}
其中，
\begin{equation}
  \pi_i=P(i_1=q_i),\qquad i=1,2,\cdots,N
  \tag{10.6}
\end{equation}
是时刻 $t=1$ 处于状态 $q_i$ 的概率。

隐马尔可夫模型由初始状态概率向量 $\pi$、状态转移概率矩阵 $A$ 和观测概率矩阵 $B$ 决定。$\pi$ 和 $A$ 决定状态序列，$B$ 决定观测序列。因此，隐马尔可夫模型 $\lambda$ 可以用三元符号表示，即
\begin{equation}
  \lambda=(A,B,\pi)
  \tag{10.7}
\end{equation}
$A,B,\pi$ 称为隐马尔可夫模型的三要素。

状态转移概率矩阵 $A$ 与初始状态概率向量 $\pi$ 确定了隐藏的马尔可夫链，生成不可观测的状态序列。观测概率矩阵 $B$ 确定了如何从状态生成观测，与状态序列综合确定了如何产生观测序列。

从定义可知，隐马尔可夫模型作了两个基本假设：

（1）齐次马尔可夫性假设，即假设隐藏的马尔可夫链在任意时刻 $t$ 的状态只依赖于其前一时刻的状态，与其他时刻的状态及观测无关，也与时刻 $t$ 无关：
\begin{equation}
  P(i_t\mid i_{t-1},o_{t-1},\cdots,i_1,o_1)=P(i_t\mid i_{t-1}),\qquad t=1,2,\cdots,T
  \tag{10.8}
\end{equation}

（2）观测独立性假设，即假设任意时刻的观测只依赖于该时刻的马尔可夫链的状态，与其他观测及状态无关：
\begin{equation}
  P(o_t\mid i_T,o_T,i_{T-1},o_{T-1},\cdots,i_{t+1},o_{t+1},i_t,i_{t-1},o_{t-1},\cdots,i_1,o_1)=P(o_t\mid i_t)
  \tag{10.9}
\end{equation}

隐马尔可夫模型可以用于标注，这时状态对应着标记。标注问题是给定观测的序列预测其对应的标记序列。可以假设标注问题的数据是由隐马尔可夫模型生成的。这样我们可以利用隐马尔可夫模型的学习与预测算法进行标注。

下面看一个隐马尔可夫模型的例子。

\noindent\textbf{例 10.1（盒子和球模型）}\quad 假设有 $4$ 个盒子，每个盒子里都装有红、白两种颜色的球，盒子里的红、白球数由表 10.1 列出。

\begin{table}[htbp]
  \centering
  \caption{各盒子的红、白球数}
  \begin{tabular}{c|cccc}
    \toprule
    盒子 & 1 & 2 & 3 & 4\\
    \midrule
    红球数 & 5 & 3 & 6 & 8\\
    白球数 & 5 & 7 & 4 & 2\\
    \bottomrule
  \end{tabular}
\end{table}

按照下面的方法抽球，产生一个球的颜色的观测序列：

\begin{itemize}
  \item 开始，从 $4$ 个盒子里以等概率随机选取 $1$ 个盒子，从这个盒子里随机抽出 $1$ 个球，记录其颜色后，放回；
  \item 然后，从当前盒子随机转移到下一个盒子，规则是：如果当前盒子是盒子 $1$，那么下一盒子一定是盒子 $2$；如果当前是盒子 $2$ 或 $3$，那么分别以概率 $0.4$ 和 $0.6$ 转移到左边或右边的盒子；如果当前是盒子 $4$，那么各以 $0.5$ 的概率停留在盒子 $4$ 或转移到盒子 $3$；
  \item 确定转移的盒子后，再从这个盒子里随机抽出 $1$ 个球，记录其颜色，放回；
  \item 如此下去，重复进行 $5$ 次，得到一个球的颜色的观测序列：
\end{itemize}
\[
  O=(\text{红},\text{红},\text{白},\text{白},\text{红}).
\]

在这个过程中，观察者只能观测到球的颜色的序列，观测不到球是从哪个盒子取出的，即观测不到盒子的序列。

在这个例子中有两个随机序列，一个是盒子的序列（状态序列），一个是球的颜色的观测序列（观测序列）。前者是隐藏的，只有后者是可观测的。这是一个隐马尔可夫模型的例子。根据所给条件，可以明确状态集合、观测集合、序列长度以及模型的三要素。

盒子对应状态，状态的集合是：
\[
  Q=\{\text{盒子 }1,\text{ 盒子 }2,\text{ 盒子 }3,\text{ 盒子 }4\},\qquad N=4.
\]
球的颜色对应观测。观测的集合是：
\[
  V=\{\text{红},\text{白}\},\qquad M=2.
\]
状态序列和观测序列长度 $T=5$。

初始概率分布为
\[
  \pi=(0.25,\,0.25,\,0.25,\,0.25)^{\mathrm T}.
\]
状态转移概率分布为
\[
  A=
  \begin{bmatrix}
    0 & 1 & 0 & 0\\
    0.4 & 0 & 0.6 & 0\\
    0 & 0.4 & 0 & 0.6\\
    0 & 0 & 0.5 & 0.5
  \end{bmatrix}.
\]
观测概率分布为
\[
  B=
  \begin{bmatrix}
    0.5 & 0.5\\
    0.3 & 0.7\\
    0.6 & 0.4\\
    0.8 & 0.2
  \end{bmatrix}.
\]

\subsection*{10.1.2\quad 观测序列的生成过程}

根据隐马尔可夫模型定义，可以将一个长度为 $T$ 的观测序列 $O=(o_1,o_2,\cdots,o_T)$ 的生成过程描述如下。

\noindent\textbf{算法 10.1（观测序列的生成）}\quad

输入：隐马尔可夫模型 $\lambda=(A,B,\pi)$，观测序列长度 $T$；

输出：观测序列 $O=(o_1,o_2,\cdots,o_T)$。

（1）按照初始状态分布 $\pi$ 产生状态 $i_1$；

（2）令 $t=1$；

（3）按照状态 $i_t$ 的观测概率分布 $b_{i_t}(k)$ 生成 $o_t$；

（4）按照状态 $i_t$ 的状态转移概率分布 $\{a_{i_t i_{t+1}}\}$ 产生状态 $i_{t+1}$，$i_{t+1}=1,2,\cdots,N$；

（5）令 $t=t+1$；如果 $t<T$，转步（3）；否则，终止。

\subsection*{10.1.3\quad 隐马尔可夫模型的 3 个基本问题}

隐马尔可夫模型有 $3$ 个基本问题：

（1）概率计算问题。给定模型 $\lambda=(A,B,\pi)$ 和观测序列 $O=(o_1,o_2,\cdots,o_T)$，计算在模型 $\lambda$ 下观测序列 $O$ 出现的概率 $P(O\mid\lambda)$。

（2）学习问题。已知观测序列 $O=(o_1,o_2,\cdots,o_T)$，估计模型 $\lambda=(A,B,\pi)$ 参数，使得在该模型下观测序列概率 $P(O\mid\lambda)$ 最大。即用极大似然估计的方法估计参数。

（3）预测问题，也称为解码（decoding）问题。已知模型 $\lambda=(A,B,\pi)$ 和观测序列 $O=(o_1,o_2,\cdots,o_T)$，求对给定观测序列条件概率 $P(I\mid O)$ 最大的状态序列 $I=(i_1,i_2,\cdots,i_T)$。即给定观测序列，求最有可能的对应状态序列。

下面各节将逐一介绍这些基本问题的解法。

\section*{10.2\quad 概率计算算法}

本节介绍计算观测序列概率 $P(O\mid\lambda)$ 的前向（forward）与后向（backward）算法。先介绍概念上可行但计算上不可行的直接计算法。

\subsection*{10.2.1\quad 直接计算法}

给定模型 $\lambda=(A,B,\pi)$ 和观测序列 $O=(o_1,o_2,\cdots,o_T)$，计算观测序列 $O$ 出现的概率 $P(O\mid\lambda)$。最直接的方法是按概率公式直接计算。通过列举所有可能的长度为 $T$ 的状态序列 $I=(i_1,i_2,\cdots,i_T)$，求各个状态序列 $I$ 与观测序列 $O=(o_1,o_2,\cdots,o_T)$ 的联合概率 $P(O,I\mid\lambda)$，然后对所有可能的状态序列求和，得到 $P(O\mid\lambda)$。

状态序列 $I=(i_1,i_2,\cdots,i_T)$ 的概率是：
\begin{equation}
  P(I\mid\lambda)=\pi_{i_1}a_{i_1i_2}a_{i_2i_3}\cdots a_{i_{T-1}i_T}
  \tag{10.10}
\end{equation}
对固定的状态序列 $I=(i_1,i_2,\cdots,i_T)$，观测序列 $O=(o_1,o_2,\cdots,o_T)$ 的概率是：
\begin{equation}
  P(O\mid I,\lambda)=b_{i_1}(o_1)b_{i_2}(o_2)\cdots b_{i_T}(o_T)
  \tag{10.11}
\end{equation}
$O$ 和 $I$ 同时出现的联合概率为
\begin{equation}
\begin{aligned}
  P(O,I\mid\lambda)
  &=P(O\mid I,\lambda)P(I\mid\lambda)\\
  &=\pi_{i_1}b_{i_1}(o_1)a_{i_1i_2}b_{i_2}(o_2)\cdots a_{i_{T-1}i_T}b_{i_T}(o_T)
\end{aligned}
  \tag{10.12}
\end{equation}
然后，对所有可能的状态序列 $I$ 求和，得到观测序列 $O$ 的概率 $P(O\mid\lambda)$，即
\begin{equation}
\begin{aligned}
  P(O\mid\lambda)
  &=\sum_I P(O\mid I,\lambda)P(I\mid\lambda)\\
  &=\sum_{i_1,i_2,\cdots,i_T}\pi_{i_1}b_{i_1}(o_1)a_{i_1i_2}b_{i_2}(o_2)\cdots a_{i_{T-1}i_T}b_{i_T}(o_T)
\end{aligned}
  \tag{10.13}
\end{equation}

但是，利用公式（10.13）计算量很大，是 $O(TN^T)$ 阶的，这种算法不可行。

下面介绍计算观测序列概率 $P(O\mid\lambda)$ 的有效算法：前向--后向算法（forward-backward algorithm）。

\subsection*{10.2.2\quad 前向算法}

首先定义前向概率。

\noindent\textbf{定义 10.2（前向概率）}\quad 给定隐马尔可夫模型 $\lambda$，定义到时刻 $t$ 部分观测序列为 $o_1,o_2,\cdots,o_t$ 且状态为 $q_i$ 的概率为前向概率，记作
\begin{equation}
  \alpha_t(i)=P(o_1,o_2,\cdots,o_t,i_t=q_i\mid\lambda)
  \tag{10.14}
\end{equation}

可以递推地求得前向概率 $\alpha_t(i)$ 及观测序列概率 $P(O\mid\lambda)$。

\noindent\textbf{算法 10.2（观测序列概率的前向算法）}\quad

输入：隐马尔可夫模型 $\lambda$，观测序列 $O$；

输出：观测序列概率 $P(O\mid\lambda)$。

（1）初值
\begin{equation}
  \alpha_1(i)=\pi_i b_i(o_1),\qquad i=1,2,\cdots,N
  \tag{10.15}
\end{equation}

（2）递推\quad 对 $t=1,2,\cdots,T-1$，
\begin{equation}
  \alpha_{t+1}(i)=\left[\sum_{j=1}^N\alpha_t(j)a_{ji}\right]b_i(o_{t+1}),\qquad i=1,2,\cdots,N
  \tag{10.16}
\end{equation}

（3）终止
\begin{equation}
  P(O\mid\lambda)=\sum_{i=1}^N\alpha_T(i)
  \tag{10.17}
\end{equation}

前向算法，步骤（1）初始化前向概率，是初始时刻的状态 $i_1=q_i$ 和观测 $o_1$ 的联合概率。步骤（2）是前向概率的递推公式，计算到时刻 $t+1$ 部分观测序列为 $o_1,o_2,\cdots,o_t,o_{t+1}$ 且在时刻 $t+1$ 处于状态 $q_i$ 的前向概率，如图 10.1 所示。在式（10.16）的方括弧里，既然 $\alpha_t(j)$ 是到时刻 $t$ 观测到 $o_1,o_2,\cdots,o_t$ 并在时刻 $t$ 处于状态 $q_j$ 的前向概率，那么乘积 $\alpha_t(j)a_{ji}$ 就是到时刻 $t$ 观测到 $o_1,o_2,\cdots,o_t$ 并在时刻 $t$ 处于状态 $q_j$ 而在时刻 $t+1$ 到达状态 $q_i$ 的联合概率。对这个乘积在时刻 $t$ 的所有可能的 $N$ 个状态 $q_j$ 求和，其结果就是到时刻 $t$ 观测为 $o_1,o_2,\cdots,o_t$ 并在时刻 $t+1$ 处于状态 $q_i$ 的联合概率。方括弧里的值与观测概率 $b_i(o_{t+1})$ 的乘积恰好是到时刻 $t+1$ 观测到 $o_1,o_2,\cdots,o_t,o_{t+1}$ 并在时刻 $t+1$ 处于状态 $q_i$ 的前向概率 $\alpha_{t+1}(i)$。步骤（3）给出 $P(O\mid\lambda)$ 的计算公式。因为
\[
  \alpha_T(i)=P(o_1,o_2,\cdots,o_T,i_T=q_i\mid\lambda)
\]
所以
\[
  P(O\mid\lambda)=\sum_{i=1}^N\alpha_T(i).
\]

\begin{figure}[htbp]
  \centering
  \includegraphics[width=0.46\textwidth]{10-1.png}
  \caption{前向概率的递推公式}
\end{figure}

如图 10.2 所示，前向算法实际是基于“状态序列的路径结构”递推计算 $P(O\mid\lambda)$ 的算法。前向算法高效的关键是其局部计算前向概率，然后利用路径结构将前向概率“递推”到全局，得到 $P(O\mid\lambda)$。具体地，在时刻 $t=1$，计算 $\alpha_1(i)$ 的 $N$ 个值（$i=1,2,\cdots,N$）；在各个时刻 $t=1,2,\cdots,T-1$，计算 $\alpha_{t+1}(i)$ 的 $N$ 个值（$i=1,2,\cdots,N$），而且每个 $\alpha_{t+1}(i)$ 的计算利用前一时刻 $N$ 个 $\alpha_t(j)$。减少计算量的原因在于每一次计算直接引用前一个时刻的计算结果，避免重复计算。这样，利用前向概率计算 $P(O\mid\lambda)$ 的计算量是 $O(N^2T)$ 阶的，而不是直接计算的 $O(TN^T)$ 阶。

\begin{figure}[htbp]
  \centering
  \includegraphics[width=0.72\textwidth]{10-2.png}
  \caption{观测序列路径结构}
\end{figure}

\noindent\textbf{例 10.2}\quad 考虑盒子和球模型 $\lambda=(A,B,\pi)$，状态集合 $Q=\{1,2,3\}$，观测集合 $V=\{\text{红},\text{白}\}$，
\[
  A=
  \begin{bmatrix}
    0.5 & 0.2 & 0.3\\
    0.3 & 0.5 & 0.2\\
    0.2 & 0.3 & 0.5
  \end{bmatrix},
  \quad
  B=
  \begin{bmatrix}
    0.5 & 0.5\\
    0.4 & 0.6\\
    0.7 & 0.3
  \end{bmatrix},
  \quad
  \pi=
  \begin{bmatrix}
    0.2\\
    0.4\\
    0.4
  \end{bmatrix}.
\]
设 $T=3$，$O=(\text{红},\text{白},\text{红})$，试用前向算法计算 $P(O\mid\lambda)$。

\noindent\textbf{解}\quad 按照算法 10.2

（1）计算初值
\[
  \alpha_1(1)=\pi_1b_1(o_1)=0.10,
\]
\[
  \alpha_1(2)=\pi_2b_2(o_1)=0.16,
\]
\[
  \alpha_1(3)=\pi_3b_3(o_1)=0.28.
\]

（2）递推计算
\[
  \alpha_2(1)=\left[\sum_{i=1}^3\alpha_1(i)a_{i1}\right]b_1(o_2)=0.154\times 0.5=0.077,
\]
\[
  \alpha_2(2)=\left[\sum_{i=1}^3\alpha_1(i)a_{i2}\right]b_2(o_2)=0.184\times 0.6=0.1104,
\]
\[
  \alpha_2(3)=\left[\sum_{i=1}^3\alpha_1(i)a_{i3}\right]b_3(o_2)=0.202\times 0.3=0.0606,
\]
\[
  \alpha_3(1)=\left[\sum_{i=1}^3\alpha_2(i)a_{i1}\right]b_1(o_3)=0.04187,
\]
\[
  \alpha_3(2)=\left[\sum_{i=1}^3\alpha_2(i)a_{i2}\right]b_2(o_3)=0.03551,
\]
\[
  \alpha_3(3)=\left[\sum_{i=1}^3\alpha_2(i)a_{i3}\right]b_3(o_3)=0.05284.
\]

（3）终止
\[
  P(O\mid\lambda)=\sum_{i=1}^3\alpha_3(i)=0.13022.
\]

\subsection*{10.2.3\quad 后向算法}

\noindent\textbf{定义 10.3（后向概率）}\quad 给定隐马尔可夫模型 $\lambda$，定义在时刻 $t$ 状态为 $q_i$ 的条件下，从 $t+1$ 到 $T$ 的部分观测序列为 $o_{t+1},o_{t+2},\cdots,o_T$ 的概率为后向概率，记作
\begin{equation}
  \beta_t(i)=P(o_{t+1},o_{t+2},\cdots,o_T\mid i_t=q_i,\lambda)
  \tag{10.18}
\end{equation}

可以用递推的方法求得后向概率 $\beta_t(i)$ 及观测序列概率 $P(O\mid\lambda)$。

\noindent\textbf{算法 10.3（观测序列概率的后向算法）}\quad

输入：隐马尔可夫模型 $\lambda$，观测序列 $O$；

输出：观测序列概率 $P(O\mid\lambda)$。

（1）
\begin{equation}
  \beta_T(i)=1,\qquad i=1,2,\cdots,N
  \tag{10.19}
\end{equation}

（2）对 $t=T-1,T-2,\cdots,1$
\begin{equation}
  \beta_t(i)=\sum_{j=1}^N a_{ij}b_j(o_{t+1})\beta_{t+1}(j),\qquad i=1,2,\cdots,N
  \tag{10.20}
\end{equation}

（3）
\begin{equation}
  P(O\mid\lambda)=\sum_{i=1}^N\pi_i b_i(o_1)\beta_1(i)
  \tag{10.21}
\end{equation}

步骤（1）初始化后向概率，对最终时刻的所有状态 $q_i$ 规定 $\beta_T(i)=1$。步骤（2）是后向概率的递推公式。如图 10.3 所示，为了计算在时刻 $t$ 状态为 $q_i$ 条件下时刻 $t+1$ 之后的观测序列为 $o_{t+1},o_{t+2},\cdots,o_T$ 的后向概率 $\beta_t(i)$，只需考虑在时刻 $t+1$ 所有可能的 $N$ 个状态 $q_j$ 的转移概率（即 $a_{ij}$ 项），以及在此状态下的观测 $o_{t+1}$ 的观测概率（即 $b_j(o_{t+1})$ 项），然后考虑状态 $q_j$ 之后的观测序列的后向概率（即 $\beta_{t+1}(j)$ 项）。步骤（3）求 $P(O\mid\lambda)$ 的思路与步骤（2）一致，只是初始概率 $\pi_i$ 代替转移概率。

\begin{figure}[htbp]
  \centering
  \includegraphics[width=0.46\textwidth]{10-3.png}
  \caption{后向概率递推公式}
\end{figure}

利用前向概率和后向概率的定义可以将观测序列概率 $P(O\mid\lambda)$ 统一写成
\begin{equation}
  P(O\mid\lambda)=\sum_{i=1}^N\sum_{j=1}^N\alpha_t(i)a_{ij}b_j(o_{t+1})\beta_{t+1}(j),\qquad t=1,2,\cdots,T-1
  \tag{10.22}
\end{equation}

\subsection*{10.2.4\quad 一些概率与期望值的计算}

利用前向概率和后向概率，可以得到关于单个状态和两个状态概率的计算公式。

1. 给定模型 $\lambda$ 和观测 $O$，在时刻 $t$ 处于状态 $q_i$ 的概率。记
\begin{equation}
  \gamma_t(i)=P(i_t=q_i\mid O,\lambda)
  \tag{10.23}
\end{equation}
可以通过前向后向概率计算。事实上，
\[
  \gamma_t(i)=P(i_t=q_i\mid O,\lambda)=\frac{P(i_t=q_i,O\mid\lambda)}{P(O\mid\lambda)}.
\]
由前向概率 $\alpha_t(i)$ 和后向概率 $\beta_t(i)$ 定义可知：
\[
  \alpha_t(i)\beta_t(i)=P(i_t=q_i,O\mid\lambda).
\]
于是得到：
\begin{equation}
  \gamma_t(i)=\frac{\alpha_t(i)\beta_t(i)}{P(O\mid\lambda)}
  =\frac{\alpha_t(i)\beta_t(i)}{\sum_{j=1}^N\alpha_t(j)\beta_t(j)}
  \tag{10.24}
\end{equation}

2. 给定模型 $\lambda$ 和观测 $O$，在时刻 $t$ 处于状态 $q_i$ 且在时刻 $t+1$ 处于状态 $q_j$ 的概率。记
\begin{equation}
  \xi_t(i,j)=P(i_t=q_i,i_{t+1}=q_j\mid O,\lambda)
  \tag{10.25}
\end{equation}
可以通过前向后向概率计算：
\[
  \xi_t(i,j)=\frac{P(i_t=q_i,i_{t+1}=q_j,O\mid\lambda)}{P(O\mid\lambda)}
  =\frac{P(i_t=q_i,i_{t+1}=q_j,O\mid\lambda)}
  {\sum_{i=1}^N\sum_{j=1}^N P(i_t=q_i,i_{t+1}=q_j,O\mid\lambda)}
\]
而
\[
  P(i_t=q_i,i_{t+1}=q_j,O\mid\lambda)=\alpha_t(i)a_{ij}b_j(o_{t+1})\beta_{t+1}(j).
\]
所以
\begin{equation}
  \xi_t(i,j)=
  \frac{\alpha_t(i)a_{ij}b_j(o_{t+1})\beta_{t+1}(j)}
  {\sum_{i=1}^N\sum_{j=1}^N\alpha_t(i)a_{ij}b_j(o_{t+1})\beta_{t+1}(j)}
  \tag{10.26}
\end{equation}

3. 将 $\gamma_t(i)$ 和 $\xi_t(i,j)$ 对各个时刻 $t$ 求和，可以得到一些有用的期望值。

（1）在观测 $O$ 下状态 $i$ 出现的期望值：
\begin{equation}
  \sum_{t=1}^T \gamma_t(i)
  \tag{10.27}
\end{equation}

（2）在观测 $O$ 下由状态 $i$ 转移的期望值：
\begin{equation}
  \sum_{t=1}^{T-1}\gamma_t(i)
  \tag{10.28}
\end{equation}

（3）在观测 $O$ 下由状态 $i$ 转移到状态 $j$ 的期望值：
\begin{equation}
  \sum_{t=1}^{T-1}\xi_t(i,j)
  \tag{10.29}
\end{equation}

\section*{10.3\quad 学习算法}

隐马尔可夫模型的学习，根据训练数据是包括观测序列和对应的状态序列还是只有观测序列，可以分别由监督学习与无监督学习实现。本节首先介绍监督学习算法，而后介绍无监督学习算法——Baum-Welch 算法（也就是 EM 算法）。

\subsection*{10.3.1\quad 监督学习方法}

假设已给训练数据包含 $S$ 个长度相同的观测序列和对应的状态序列
\[
  \{(O_1,I_1),(O_2,I_2),\cdots,(O_S,I_S)\},
\]
那么可以利用极大似然估计法来估计隐马尔可夫模型的参数。具体方法如下。

1. 转移概率 $a_{ij}$ 的估计

设样本中时刻 $t$ 处于状态 $i$ 时刻 $t+1$ 转移到状态 $j$ 的频数为 $A_{ij}$，那么状态转移概率 $a_{ij}$ 的估计是
\begin{equation}
  \hat a_{ij}=\frac{A_{ij}}{\sum_{j=1}^N A_{ij}},\qquad i=1,2,\cdots,N;\quad j=1,2,\cdots,N
  \tag{10.30}
\end{equation}

2. 观测概率 $b_j(k)$ 的估计

设样本中状态为 $j$ 并观测为 $k$ 的频数是 $B_{jk}$，那么状态为 $j$ 观测为 $k$ 的概率 $b_j(k)$ 的估计是
\begin{equation}
  \hat b_j(k)=\frac{B_{jk}}{\sum_{k=1}^M B_{jk}},\qquad j=1,2,\cdots,N;\quad k=1,2,\cdots,M
  \tag{10.31}
\end{equation}

3. 初始状态概率 $\pi_i$ 的估计 $\hat\pi_i$ 为 $S$ 个样本中初始状态为 $q_i$ 的频率

由于监督学习需要使用标注的训练数据，而人工标注训练数据往往代价很高，有时就会利用无监督学习的方法。

\subsection*{10.3.2\quad Baum-Welch 算法}

假设给定训练数据只包含 $S$ 个长度为 $T$ 的观测序列 $\{O_1,O_2,\cdots,O_S\}$ 而没有对应的状态序列，目标是学习隐马尔可夫模型 $\lambda=(A,B,\pi)$ 的参数。我们将观测序列数据看作观测数据 $O$，状态序列数据看作不可观测的隐数据 $I$，那么隐马尔可夫模型事实上是一个含有隐变量的概率模型
\begin{equation}
  P(O\mid\lambda)=\sum_I P(O\mid I,\lambda)P(I\mid\lambda)
  \tag{10.32}
\end{equation}
它的参数学习可以由 EM 算法实现。

1. 确定完全数据的对数似然函数

所有观测数据写成 $O=(o_1,o_2,\cdots,o_T)$，所有隐数据写成 $I=(i_1,i_2,\cdots,i_T)$，完全数据是 $(O,I)=(o_1,o_2,\cdots,o_T,i_1,i_2,\cdots,i_T)$。完全数据的对数似然函数是 $\log P(O,I\mid\lambda)$。

2. EM 算法的 E 步：求 $Q$ 函数 $Q(\lambda,\bar\lambda)$\footnote{按照 $Q$ 函数的定义
\[
  Q(\lambda,\bar\lambda)=E_I[\log P(O,I\mid\lambda)\mid O,\bar\lambda]
\]
式（10.33）略去了对 $\lambda$ 而言的常数因子 $1/P(O\mid\lambda)$。}
\begin{equation}
  Q(\lambda,\bar\lambda)=\sum_I \log P(O,I\mid\lambda)P(O,I\mid\bar\lambda)
  \tag{10.33}
\end{equation}
其中，$\bar\lambda$ 是隐马尔可夫模型参数的当前估计值，$\lambda$ 是要极大化的隐马尔可夫模型参数。
\[
  P(O,I\mid\lambda)=\pi_{i_1}b_{i_1}(o_1)a_{i_1i_2}b_{i_2}(o_2)\cdots a_{i_{T-1}i_T}b_{i_T}(o_T)
\]
于是函数 $Q(\lambda,\bar\lambda)$ 可以写成：
\begin{equation}
\begin{aligned}
  Q(\lambda,\bar\lambda)
  &=
  \sum_I \log\pi_{i_1}P(O,I\mid\bar\lambda)
  +\sum_I\left(\sum_{t=1}^{T-1}\log a_{i_t i_{t+1}}\right)P(O,I\mid\bar\lambda)\\
  &\quad
  +\sum_I\left(\sum_{t=1}^{T}\log b_{i_t}(o_t)\right)P(O,I\mid\bar\lambda)
\end{aligned}
  \tag{10.34}
\end{equation}
式中求和都是对所有数据的序列总长度 $T$ 进行的。

3. EM 算法的 M 步：极大化 $Q$ 函数 $Q(\lambda,\bar\lambda)$ 求模型参数 $A,B,\pi$

由于要极大化的参数在式（10.34）中单独地出现在 $3$ 个项中，所以只需对各项分别极大化。

（1）式（10.34）的第 $1$ 项可以写成：
\[
  \sum_I \log\pi_{i_1}P(O,I\mid\bar\lambda)=\sum_{i=1}^N\log\pi_i P(O,i_1=i\mid\bar\lambda)
\]
注意到 $\pi_i$ 满足约束条件 $\sum_{i=1}^N\pi_i=1$，利用拉格朗日乘子法，写出拉格朗日函数：
\[
  \sum_{i=1}^N\log\pi_i P(O,i_1=i\mid\bar\lambda)
  +\gamma\left(\sum_{i=1}^N\pi_i-1\right)
\]
对其求偏导数并令结果为 $0$
\begin{equation}
  \frac{\partial}{\partial\pi_i}
  \left[
  \sum_{i=1}^N\log\pi_i P(O,i_1=i\mid\bar\lambda)
  +\gamma\left(\sum_{i=1}^N\pi_i-1\right)
  \right]=0
  \tag{10.35}
\end{equation}
得
\[
  P(O,i_1=i\mid\bar\lambda)+\gamma\pi_i=0
\]
对 $i$ 求和得到 $\gamma$
\[
  \gamma=-P(O\mid\bar\lambda)
\]
代入式（10.35）即得
\begin{equation}
  \pi_i=\frac{P(O,i_1=i\mid\bar\lambda)}{P(O\mid\bar\lambda)}
  \tag{10.36}
\end{equation}

（2）式（10.34）的第 $2$ 项可以写成
\[
  \sum_I\left(\sum_{t=1}^{T-1}\log a_{i_t i_{t+1}}\right)P(O,I\mid\bar\lambda)
  =
  \sum_{i=1}^N\sum_{j=1}^N\sum_{t=1}^{T-1}\log a_{ij}P(O,i_t=i,i_{t+1}=j\mid\bar\lambda)
\]
类似第 $1$ 项，应用具有约束条件 $\sum_{j=1}^N a_{ij}=1$ 的拉格朗日乘子法可以求出
\begin{equation}
  a_{ij}=
  \frac{\sum_{t=1}^{T-1}P(O,i_t=i,i_{t+1}=j\mid\bar\lambda)}
  {\sum_{t=1}^{T-1}P(O,i_t=i\mid\bar\lambda)}
  \tag{10.37}
\end{equation}

（3）式（10.34）的第 $3$ 项为
\[
  \sum_I\left(\sum_{t=1}^{T}\log b_{i_t}(o_t)\right)P(O,I\mid\bar\lambda)
  =
  \sum_{j=1}^N\sum_{t=1}^{T}\log b_j(o_t)P(O,i_t=j\mid\bar\lambda)
\]
同样用拉格朗日乘子法，约束条件是 $\sum_{k=1}^M b_j(k)=1$。注意，只有在 $o_t=v_k$ 时 $b_j(o_t)$ 对 $b_j(k)$ 的偏导数才不为 $0$，以 $I(o_t=v_k)$ 表示。求得
\begin{equation}
  b_j(k)=
  \frac{\sum_{t=1}^{T}P(O,i_t=j\mid\bar\lambda)I(o_t=v_k)}
  {\sum_{t=1}^{T}P(O,i_t=j\mid\bar\lambda)}
  \tag{10.38}
\end{equation}

\subsection*{10.3.3\quad Baum-Welch 模型参数估计公式}

将式（10.36）～式（10.38）中的各概率分别用 $\gamma_t(i)$，$\xi_t(i,j)$ 表示，则可将相应的公式写成：
\begin{equation}
  a_{ij}=\frac{\sum_{t=1}^{T-1}\xi_t(i,j)}{\sum_{t=1}^{T-1}\gamma_t(i)}
  \tag{10.39}
\end{equation}
\begin{equation}
  b_j(k)=\frac{\sum_{t=1,o_t=v_k}^{T}\gamma_t(j)}{\sum_{t=1}^{T}\gamma_t(j)}
  \tag{10.40}
\end{equation}
\begin{equation}
  \pi_i=\gamma_1(i)
  \tag{10.41}
\end{equation}
其中，$\gamma_t(i)$，$\xi_t(i,j)$ 分别由式（10.24）及式（10.26）给出。式（10.39）～式（10.41）就是 Baum-Welch 算法（Baum-Welch algorithm），它是 EM 算法在隐马尔可夫模型学习中的具体实现，由 Baum 和 Welch 提出。

\noindent\textbf{算法 10.4（Baum-Welch 算法）}\quad

输入：观测数据 $O=(o_1,o_2,\cdots,o_T)$；

输出：隐马尔可夫模型参数。

（1）初始化。对 $n=0$，选取 $a_{ij}^{(0)}$，$b_j(k)^{(0)}$，$\pi_i^{(0)}$，得到模型 $\lambda^{(0)}=(A^{(0)},B^{(0)},\pi^{(0)})$。

（2）递推。对 $n=1,2,\cdots$，
\[
  a_{ij}^{(n+1)}
  =\frac{\sum_{t=1}^{T-1}\xi_t(i,j)}{\sum_{t=1}^{T-1}\gamma_t(i)}
\]
\[
  b_j(k)^{(n+1)}
  =\frac{\sum_{t=1,o_t=v_k}^{T}\gamma_t(j)}{\sum_{t=1}^{T}\gamma_t(j)}
\]
\[
  \pi_i^{(n+1)}=\gamma_1(i)
\]
右端各值按观测 $O=(o_1,o_2,\cdots,o_T)$ 和模型 $\lambda^{(n)}=(A^{(n)},B^{(n)},\pi^{(n)})$ 计算。式中 $\gamma_t(i)$，$\xi_t(i,j)$ 由式（10.24）和式（10.26）给出。

（3）终止。得到模型参数 $\lambda^{(n+1)}=(A^{(n+1)},B^{(n+1)},\pi^{(n+1)})$。

\section*{10.4\quad 预测算法}

下面介绍隐马尔可夫模型预测的两种算法：近似算法与维特比算法（Viterbi algorithm）。

\subsection*{10.4.1\quad 近似算法}

近似算法的想法是，在每个时刻 $t$ 选择在该时刻最有可能出现的状态 $i_t^*$，从而得到一个状态序列 $I^*=(i_1^*,i_2^*,\cdots,i_T^*)$，将它作为预测的结果。

给定隐马尔可夫模型 $\lambda$ 和观测序列 $O$，在时刻 $t$ 处于状态 $q_i$ 的概率 $\gamma_t(i)$ 是
\begin{equation}
  \gamma_t(i)=\frac{\alpha_t(i)\beta_t(i)}{P(O\mid\lambda)}
  =\frac{\alpha_t(i)\beta_t(i)}{\sum_{j=1}^N\alpha_t(j)\beta_t(j)}
  \tag{10.42}
\end{equation}
在每一时刻 $t$ 最有可能的状态 $i_t^*$ 是
\begin{equation}
  i_t^*=\arg\max_{1\le i\le N}[\gamma_t(i)],\qquad t=1,2,\cdots,T
  \tag{10.43}
\end{equation}
从而得到状态序列 $I^*=(i_1^*,i_2^*,\cdots,i_T^*)$。

近似算法的优点是计算简单，其缺点是不能保证预测的状态序列整体是最有可能的状态序列，因为预测的状态序列可能有实际不发生的部分。事实上，上述方法得到的状态序列中有可能存在转移概率为 $0$ 的相邻状态，即对某些 $i,j$，$a_{ij}=0$ 时。尽管如此，近似算法仍然是有用的。

\subsection*{10.4.2\quad 维特比算法}

维特比算法实际是用动态规划（dynamic programming）解隐马尔可夫模型预测问题，即用动态规划求概率最大路径（最优路径）。这时一条路径对应着一个状态序列。

根据动态规划原理，最优路径具有这样的特性：如果最优路径在时刻 $t$ 通过结点 $i_t^*$，那么这一条路径从结点 $i_t^*$ 到终点 $i_T^*$ 的部分路径，对于从 $i_t^*$ 到 $i_T^*$ 的所有可能的部分路径来说，必须是最优的。因为假如不是这样，那么从 $i_t^*$ 到 $i_T^*$ 就有另一条更好的部分路径存在，如果把它和从 $i_1^*$ 到达 $i_t^*$ 的部分路径连接起来，就会形成一条比原来的路径更优的路径，这是矛盾的。依据这一原理，我们只需从时刻 $t=1$ 开始，递推地计算在时刻 $t$ 状态为 $i$ 的各条部分路径的最大概率，直至得到时刻 $t=T$ 状态为 $i$ 的各条路径的最大概率。时刻 $t=T$ 的最大概率即为最优路径的概率 $P^*$，最优路径的终结点 $i_T^*$ 也同时得到。之后，为了找出最优路径的各个结点，从终结点 $i_T^*$ 开始，由后向前逐步求得结点 $i_{T-1}^*,\cdots,i_1^*$，得到最优路径 $I^*=(i_1^*,i_2^*,\cdots,i_T^*)$。这就是维特比算法。

首先导入两个变量 $\delta$ 和 $\Psi$。定义在时刻 $t$ 状态为 $i$ 的所有单个路径 $(i_1,i_2,\cdots,i_t)$ 中概率最大值为
\begin{equation}
  \delta_t(i)=\max_{i_1,i_2,\cdots,i_{t-1}}P(i_t=i,i_{t-1},\cdots,i_1,o_t,\cdots,o_1\mid\lambda),
  \qquad i=1,2,\cdots,N
  \tag{10.44}
\end{equation}
由定义可得变量 $\delta$ 的递推公式：
\begin{equation}
\begin{aligned}
  \delta_{t+1}(i)
  &=\max_{i_1,i_2,\cdots,i_t}P(i_{t+1}=i,i_t,\cdots,i_1,o_{t+1},\cdots,o_1\mid\lambda)\\
  &=\max_{1\le j\le N}[\delta_t(j)a_{ji}]b_i(o_{t+1}),\qquad
  i=1,2,\cdots,N;\quad t=1,2,\cdots,T-1
\end{aligned}
  \tag{10.45}
\end{equation}

定义在时刻 $t$ 状态为 $i$ 的所有单个路径 $(i_1,i_2,\cdots,i_{t-1},i)$ 中概率最大的路径的第 $t-1$ 个结点为
\begin{equation}
  \Psi_t(i)=\arg\max_{1\le j\le N}[\delta_{t-1}(j)a_{ji}],\qquad i=1,2,\cdots,N
  \tag{10.46}
\end{equation}

下面介绍维特比算法。

\noindent\textbf{算法 10.5（维特比算法）}\quad

输入：模型 $\lambda=(A,B,\pi)$ 和观测 $O=(o_1,o_2,\cdots,o_T)$；

输出：最优路径 $I^*=(i_1^*,i_2^*,\cdots,i_T^*)$。

（1）初始化
\[
  \delta_1(i)=\pi_i b_i(o_1),\qquad i=1,2,\cdots,N
\]
\[
  \Psi_1(i)=0,\qquad i=1,2,\cdots,N
\]

（2）递推。对 $t=2,3,\cdots,T$
\[
  \delta_t(i)=\max_{1\le j\le N}[\delta_{t-1}(j)a_{ji}]b_i(o_t),\qquad i=1,2,\cdots,N
\]
\[
  \Psi_t(i)=\arg\max_{1\le j\le N}[\delta_{t-1}(j)a_{ji}],\qquad i=1,2,\cdots,N
\]

（3）终止
\[
  P^*=\max_{1\le i\le N}\delta_T(i)
\]
\[
  i_T^*=\arg\max_{1\le i\le N}[\delta_T(i)]
\]

（4）最优路径回溯。对 $t=T-1,T-2,\cdots,1$
\[
  i_t^*=\Psi_{t+1}(i_{t+1}^*)
\]
求得最优路径 $I^*=(i_1^*,i_2^*,\cdots,i_T^*)$。

下面通过一个例子来说明维特比算法。

\noindent\textbf{例 10.3}\quad 例 10.2 的模型 $\lambda=(A,B,\pi)$，
\[
  A=
  \begin{bmatrix}
    0.5 & 0.2 & 0.3\\
    0.3 & 0.5 & 0.2\\
    0.2 & 0.3 & 0.5
  \end{bmatrix},
  \quad
  B=
  \begin{bmatrix}
    0.5 & 0.5\\
    0.4 & 0.6\\
    0.7 & 0.3
  \end{bmatrix},
  \quad
  \pi=
  \begin{bmatrix}
    0.2\\
    0.4\\
    0.4
  \end{bmatrix}.
\]
已知观测序列 $O=(\text{红},\text{白},\text{红})$，试求最优状态序列，即最优路径 $I^*=(i_1^*,i_2^*,i_3^*)$。

\noindent\textbf{解}\quad 如图 10.4 所示，要在所有可能的路径中选择一条最优路径，按照以下步骤处理：

（1）初始化。在 $t=1$ 时，对每一个状态 $i$，$i=1,2,3$，求状态为 $i$ 观测 $o_1$ 为红的概率，记此概率为 $\delta_1(i)$，则
\[
  \delta_1(i)=\pi_i b_i(o_1)=\pi_i b_i(\text{红}),\qquad i=1,2,3
\]
代入实际数据
\[
  \delta_1(1)=0.10,\qquad \delta_1(2)=0.16,\qquad \delta_1(3)=0.28.
\]
记 $\Psi_1(i)=0$，$i=1,2,3$。

\begin{figure}[htbp]
  \centering
  \includegraphics[width=0.68\textwidth]{10-4.png}
  \caption{求最优路径}
\end{figure}

（2）在 $t=2$ 时，对每个状态 $i$，$i=1,2,3$，求在 $t=1$ 时状态为 $j$ 观测为红并在 $t=2$ 时状态为 $i$ 观测 $o_2$ 为白的路径的最大概率，记此最大概率为 $\delta_2(i)$，则
\[
  \delta_2(i)=\max_{1\le j\le3}[\delta_1(j)a_{ji}]b_i(o_2)
\]
同时，对每个状态 $i$，$i=1,2,3$，记录概率最大路径的前一个状态 $j$：
\[
  \Psi_2(i)=\arg\max_{1\le j\le3}[\delta_1(j)a_{ji}],\qquad i=1,2,3
\]
计算：
\[
\begin{aligned}
  \delta_2(1)
  &=\max_{1\le j\le3}[\delta_1(j)a_{j1}]b_1(o_2)\\
  &=\max_j\{0.10\times0.5,0.16\times0.3,0.28\times0.2\}\times0.5\\
  &=0.028
\end{aligned}
\]
\[
  \Psi_2(1)=3
\]
\[
  \delta_2(2)=0.0504
\]
\[
  \Psi_2(2)=3
\]
\[
  \delta_2(3)=0.042
\]
\[
  \Psi_2(3)=3
\]

同样，在 $t=3$ 时，
\[
  \delta_3(i)=\max_{1\le j\le3}[\delta_2(j)a_{ji}]b_i(o_3)
\]
\[
  \Psi_3(i)=\arg\max_{1\le j\le3}[\delta_2(j)a_{ji}]
\]
\[
  \delta_3(1)=0.00756
\]
\[
  \Psi_3(1)=2
\]
\[
  \delta_3(2)=0.01008
\]
\[
  \Psi_3(2)=2
\]
\[
  \delta_3(3)=0.0147
\]
\[
  \Psi_3(3)=3
\]

（3）以 $P^*$ 表示最优路径的概率，则
\[
  P^*=\max_{1\le i\le3}\delta_3(i)=0.0147
\]
最优路径的终点是 $i_3^*$：
\[
  i_3^*=\arg\max_i[\delta_3(i)]=3.
\]

（4）由最优路径的终点 $i_3^*$，逆向找到 $i_2^*,i_1^*$：
\[
  \text{在 }t=2\text{ 时，}\qquad i_2^*=\Psi_3(i_3^*)=\Psi_3(3)=3
\]
\[
  \text{在 }t=1\text{ 时，}\qquad i_1^*=\Psi_2(i_2^*)=\Psi_2(3)=3
\]
于是求得最优路径，即最优状态序列 $I^*=(i_1^*,i_2^*,i_3^*)=(3,3,3)$。

\section*{本章概要}

1. 隐马尔可夫模型是关于时序的概率模型，描述由一个隐藏的马尔可夫链随机生成不可观测的状态的序列，再由各个状态随机生成一个观测从而产生观测序列的过程。

隐马尔可夫模型由初始状态概率向量 $\pi$、状态转移概率矩阵 $A$ 和观测概率矩阵 $B$ 决定。因此，隐马尔可夫模型可以写成 $\lambda=(A,B,\pi)$。

隐马尔可夫模型是一个生成模型，表示状态序列和观测序列的联合分布，但是状态序列是隐藏的，不可观测的。

隐马尔可夫模型可以用于标注，这时状态对应着标记。标注问题是给定观测序列预测其对应的标记序列。

2. 概率计算问题。给定模型 $\lambda=(A,B,\pi)$ 和观测序列 $O=(o_1,o_2,\cdots,o_T)$，计算在模型 $\lambda$ 下观测序列 $O$ 出现的概率 $P(O\mid\lambda)$。前向--后向算法通过递推地计算前向--后向概率可以高效地进行隐马尔可夫模型的概率计算。

3. 学习问题。已知观测序列 $O=(o_1,o_2,\cdots,o_T)$，估计模型 $\lambda=(A,B,\pi)$ 参数，使得在该模型下观测序列概率 $P(O\mid\lambda)$ 最大。即用极大似然估计的方法估计参数。Baum-Welch 算法，也就是 EM 算法可以高效地对隐马尔可夫模型进行训练。它是一种无监督学习算法。

4. 预测问题。已知模型 $\lambda=(A,B,\pi)$ 和观测序列 $O=(o_1,o_2,\cdots,o_T)$，求对给定观测序列条件概率 $P(I\mid O)$ 最大的状态序列 $I=(i_1,i_2,\cdots,i_T)$。维特比算法应用动态规划高效地求解最优路径，即概率最大的状态序列。

\section*{继续阅读}

隐马尔可夫模型的介绍可见文献 [1,2]，特别地，文献 [1] 是经典的介绍性论文。关于 Baum-Welch 算法可见文献 [3,4]。可以认为概率上下文无关文法（probabilistic context-free grammar）是隐马尔可夫模型的一种推广，隐马尔可夫模型的不可观测数
据是状态序列，而概率上下文无关文法的不可观测数据是上下文无关文法树 [5]。动态贝叶斯网络（dynamic Bayesian network）是定义在时序数据上的贝叶斯网络，它包含隐马尔可夫模型，是一种特例 [6]。

\section*{习题}

\noindent\textbf{10.1}\quad 给定盒子和球组成的隐马尔可夫模型 $\lambda=(A,B,\pi)$，其中，
\[
  A=
  \begin{bmatrix}
    0.5 & 0.2 & 0.3\\
    0.3 & 0.5 & 0.2\\
    0.2 & 0.3 & 0.5
  \end{bmatrix},
  \quad
  B=
  \begin{bmatrix}
    0.5 & 0.5\\
    0.4 & 0.6\\
    0.7 & 0.3
  \end{bmatrix},
  \quad
  \pi=(0.2,0.4,0.4)^{\mathrm T}
\]
设 $T=4$，$O=(\text{红},\text{白},\text{红},\text{白})$，试用后向算法计算 $P(O\mid\lambda)$。

\noindent\textbf{10.2}\quad 考虑盒子和球组成的隐马尔可夫模型 $\lambda=(A,B,\pi)$，其中，
\[
  A=
  \begin{bmatrix}
    0.5 & 0.1 & 0.4\\
    0.3 & 0.5 & 0.2\\
    0.2 & 0.2 & 0.6
  \end{bmatrix},
  \quad
  B=
  \begin{bmatrix}
    0.5 & 0.5\\
    0.4 & 0.6\\
    0.7 & 0.3
  \end{bmatrix},
  \quad
  \pi=(0.2,0.3,0.5)^{\mathrm T}
\]
设 $T=8$，$O=(\text{红},\text{白},\text{红},\text{红},\text{白},\text{红},\text{白},\text{白})$，用前向后向概率计算 $P(i_4=q_3\mid O,\lambda)$。

\noindent\textbf{10.3}\quad 在习题 10.1 中，试用维特比算法求最优路径 $I^*=(i_1^*,i_2^*,i_3^*,i_4^*)$。

\noindent\textbf{10.4}\quad 试用前向概率和后向概率推导
\[
  P(O\mid\lambda)=\sum_{i=1}^N\sum_{j=1}^N\alpha_t(i)a_{ij}b_j(o_{t+1})\beta_{t+1}(j),
  \qquad t=1,2,\cdots,T-1
\]

\noindent\textbf{10.5}\quad 比较维特比算法中变量 $\delta$ 的计算和前向算法中变量 $\alpha$ 的计算的主要区别。

\section*{参考文献}

\noindent [1] Rabiner L, Juang B. An introduction to hidden Markov Models. IEEE ASSP Magazine, 1986, 3(1): 4--16.

\noindent [2] Rabiner L. A tutorial on hidden Markov models and selected applications in speech recognition. Proceedings of IEEE, 1989, 77(2): 257--286.

\noindent [3] Baum L, et al. A maximization technique occurring in the statistical analysis of probabilistic functions of Markov chains. Annals of Mathematical Statistics, 1970, 41: 164--171.

\noindent [4] Bilmes J A. A gentle tutorial of the EM algorithm and its application to parameter estimation for Gaussian mixture and hidden Markov models. \url{http://ssli.ee.washington.edu/~bilmes/mypubs/bilmes1997-em.pdf}.

\noindent [5] Lari K, Young S J. Applications of stochastic context-free grammars using the Inside-Outside algorithm. Computer Speech \& Language, 1991, 5(3): 237--257.

\noindent [6] Ghahramani Z. Learning dynamic Bayesian networks. Lecture Notes in Computer Science, Vol. 1387, Springer, 1997, 168--197.

\end{document}
